\chapter{Introduction}
\label{chap:introduction}
High-order numerical methods represent smooth solutions accurately with
relatively few degrees of freedom. Conservation laws also produce shocks,
where that same high-order representation can oscillate. Discontinuous Galerkin
(DG) combines finite element approximation inside each mesh cell with numerical
fluxes between cells. The central practical question here is how to control
unresolved oscillations while retaining useful smooth variation.

\section{Aim and contribution}
This thesis develops a reusable MFEM realization of a specifically defined
adapted oscillation-free discontinuous Galerkin (OFDG) method. It combines projection damping with face-averaged jump
sensing and optional KXRCF activation. The implementation operates directly in
the finite element coefficient basis, including nodal bases, and extends the
construction to static polynomial curved meshes. Repeated physical derivatives
on these meshes are approximated by recursive projection, with tests measuring
their accuracy and geometric consistency.

The contribution has three connected parts:
\begin{enumerate}
 \item a clear definition and maintainable implementation of the adapted filter,
 including the distinctions from published OFDG and OEDG formulations;
 \item verification of conservation, basis independence, geometric operators,
 and serial/two-rank consistency within an explicit support boundary;
 \item preliminary numerical characterization of damping and KXRCF gating on
 affordable, physically meaningful examples.
\end{enumerate}
The implementation is held in the project repository
\url{https://github.com/tzuend/OFDG}. The report distinguishes algebraic
identities, numerical verification, and empirical observations. These support
different kinds of statements and should not be substituted for one another.

\section{Related work and focused questions}
Projection-based oscillation damping was developed for scalar conservation laws
by Lu, Liu, and Shu~\cite{LuEtAl2021} and extended to systems
in~\cite{LiuEtAl2022}. Peng, Sun, and Wu introduce oscillation-eliminating DG (OEDG), with
normalized damping and stage-wise exact attenuation~\cite{PengEtAl2024}.
KXRCF supplies a separate troubled-cell decision, motivated by accurate DG
outflow traces~\cite{KrivodonovaEtAl2004}. Detector comparisons show sensitivity
to variable choice, order, and normalization~\cite{QiuShu2005,FuShu2017}.

This project asks: can these mechanisms be realized without imposing a modal
storage basis; can their essential invariants survive curved geometry and MPI
partitioning; and what does selective damping change in the measured solutions?
The final question remains only partly answered. A detector-activity study is
still needed to explain where gating helps or fails. Performance attribution
and a general robustness ranking require evidence beyond the current preview.

\section{How to read this working report}
\Cref{chap:methods} bridges from familiar finite element ideas to DG traces,
fluxes, and time stepping. \Cref{chap:stabilization} defines adapted OFDG,
KXRCF, OEDG, and positivity treatment. \Cref{chap:implementation} describes
the support boundary and verification. \Cref{chap:experiments} presents the
measured preview, followed by interpretation in \cref{chap:discussion} and
conclusions in \cref{chap:conclusion}. Detailed matrix construction and
reproducibility instructions are in the appendices.

\EvidenceBox{EVIDENCE AND COMPLETION STATUS}{A \textbf{preliminary result} is
an actual measurement with stated limits. A \textbf{PLACEHOLDER} lists specific
missing evidence for an intended stronger argument. A \textbf{scope limitation}
is an explicit exclusion, not a promised future result. The register in
\cref{app:missing-work} collects unresolved work. No mock numerical results or
filler discussion are used.}
